\section{Mise en œuvre et exploitation}\label{sec-ops}

\subsection{Package}
Le package Python \code{repo\_control} (numpy et scipy ; scikit-learn uniquement pour la référence par mélange gaussien)
implémente tous les objets de cette note :
\begin{center}\small
\begin{tabular}{p{2.9cm}p{12.2cm}}
\toprule
Module & Contenu\\
\midrule
\code{curves} & courbes zéro-coupon à forwards constants, dividendes discrets, forwards, bootstrap de l'\cref{alg-bootstrap}\\
\code{parity} & forward et facteur d'actualisation à partir des prix de consensus des calls et des puts (\cref{lem-parity})\\
\code{buckets} & périodes de dividendes, allocation, matrice d'allocation $A$ (\cref{lem-allocation})\\
\code{noise} & écart de panel, covariance de panels emboîtés, covariances des panels options et dividendes, plafonnement de la dispersion (\cref{sec-noise})\\
\code{pipeline} & instantanés Totem et FO $\to$ repo implicite, écart de portage, $\Sigma$, contexte de détection\\
\code{detect} & gabarits, filtres adaptés, calibrations familiale et de $\kappa$, attribution, réparation, test d'erreur de marquage FO\\
\code{control} & lignes de contrôle, jacobienne du bootstrap $J$, matérialité de la répartition (\cref{sec-control})\\
\code{pricer} & un pricer minimal qui produit les sensibilités par période pour les exemples\\
\code{synthetic}, \code{benchmark}, \code{baselines} & marché synthétique construit contributeur par contributeur et banc d'essai de la \cref{sec-benchmark}\\
\bottomrule
\end{tabular}
\end{center}
Chaque composant a un contrôle en forme fermée dans la suite de tests : aller-retour du bootstrap, parité, identité de
contrôle contre la réévaluation complète, formules d'écart de panel et de sortie de contributeurs, calibration du
détecteur sous l'hypothèse nulle et reconnaissance de chaque gabarit.

\subsection{Données de production}
\begin{center}\small
\begin{tabular}{p{2.4cm}p{10.0cm}p{2.4cm}}
\toprule
Donnée & Champs & Forme\\
\midrule
Options Totem & échéances ; strikes ; prix de consensus des calls et des puts ; contributeurs ; dispersion des forwards des contributeurs (ou une estimation) ; spot de référence & $K\times M$\\
Dividendes Totem & bornes des périodes ; niveau de consensus ; dispersion ; contributeurs (net des dividendes déjà versés) & $Y$\\
Marché FO & spot à l'heure de l'instantané Totem ; courbe des taux ; courbe de repo zéro-coupon aux échéances Totem ; prévision de dividendes & $K$ ; liste\\
Risque FO & sensibilités repo par pilier (projetées sur les échéances Totem) ; sensibilités dividendes par période Totem & $K+Y$ par portefeuille\\
Historique & douze à vingt-quatre instantanés mensuels et incidents connus & pour $\kappa$ et la validation\\
\bottomrule
\end{tabular}
\end{center}
Si la dispersion des forwards des contributeurs n'est pas publiée, on peut l'estimer à partir de la dispersion des calls et
des puts très dans la monnaie, ou la calibrer sur l'historique ; le détecteur a besoin de la bonne échelle par échéance, et
$\kappa$ corrige un biais commun.

\subsection{Procédure mensuelle}
\begin{enumerate}
\item Charger les instantanés Totem et FO à la même heure de valorisation ; vérifier le spot de référence.
\item Impliciter les forwards par parité ; examiner le facteur d'actualisation implicite, les résidus de parité et les
rapports de dispersion.
\item Allouer les périodes de dividendes Totem sur le profil de détachement FO ; impliciter le repo avec les dividendes
Totem et avec les dividendes FO.
\item Construire $\Sigma$ et lancer le détecteur avec le $\kappa$ et l'échelle de niveau de la région et de la catégorie.
\item Calculer les lignes de contrôle par portefeuille : lignes forward, dividende, non-forward et repo ; vérifier la
matérialité de la répartition.
\item Examiner les alertes avec leurs explications et les erreurs de marquage FO possibles ; utiliser la courbe réparée
lorsqu'une anomalie Totem est confirmée.
\item Consigner la décision prise pour chaque alerte ; ces étiquettes alimentent la calibration suivante.
\end{enumerate}

\subsection{Calibration et gouvernance}
$\kappa$ et l'échelle de niveau sont recalibrés chaque trimestre par région et par catégorie sur un historique sain, avec
un budget d'alertes de \vFPR{} des sous-jacents par mois. Le taux d'alerte réalisé et la part des explications acceptées
par les relecteurs sont suivis. Là où ils sont cotés, des ancres de marché complètent le détecteur : contrats à terme
cotés sur indices et sur dividendes, contrats à terme de rendement total qui cotent un spread de financement, et
commissions de prêt-emprunt de titres pour le court terme des actions.

\section{Conclusion}
Les courbes de repo implicite sont irrégulières parce qu'elles sont le résidu de deux moyennes de consensus calculées sur
des panels différents ; une partie de ce résidu ne peut être supprimée par aucune méthode qui ne voit que les données de
consensus (\cref{thm-nonid}). Deux conséquences en découlent. Pour une exposition linéaire en forward, le contrôle total
n'a pas besoin de la répartition, exactement en réévaluation complète et au premier ordre sur les lignes de contrôle
(\cref{thm-control,thm-exact}), et l'exposition non-forward $g$ indique, portefeuille par portefeuille, l'importance de
la répartition. Pour la détection, les mécanismes qui produisent réellement des courbes aberrantes ont des signatures
connues dans l'espace de portage ; des filtres adaptés, avec un modèle de bruit construit à partir des propres statistiques
de Totem, les détectent à un taux de fausses alertes calibré, les expliquent en unités naturelles et les distinguent des
véritables erreurs de marquage FO. L'étape suivante consiste à faire tourner la chaîne sur douze à vingt-quatre mois
d'instantanés Totem et FO réels, à calibrer les réglages par région et par catégorie, et à mesurer l'attribution sur des
incidents étiquetés.
